% ---
% id: optics-2026-004
% title: The Great Gig in the Sky
% title_en: The Great Gig in the Sky
% topic: observations
% subtopic: ccd-photometry
% difficulty: 4
% source: camp
% event: Оптика I
% year: 2026
% grade: 9-11
% lang: ru
% ---
Gravitational-wave Optical Transient Observer предназначен для поиска
оптических остатков источников гравитационных волн различной природы
в течение одного часа после их обнаружения системами LIGO и VIRGO.
Система GOTO состоит из 4 одинаковых рефлекторов с апертурой
$\qty{40}{см}$ и относительным отверстием $1/2.5$, одновременно
фотографирующих крупный участок неба. Для простоты положим, что поля
телескопов не перекрываются.

\begin{enumerate}
  \item Вычислите масштаб изображений, получаемых на этих телескопах
        в их фокальных плоскостях, в угловых единицах на мм.
\end{enumerate}

Примем, что основными источниками ошибок матрицы являются шум считывания
$\mathrm{RON} = 10$ отсчётов на пиксел и тепловой шум $\mathrm{DN}$,
составляющий 1 отсчёт в минуту на пиксел. Используемые ПЗС-матрицы имеют
пикселы размером $\qty{6}{мкм}$ и усиление $\mathrm{gain} = 1$. Типичный размер
атмосферных изображений звёзд $\mathrm{seeing} = \asec{1.0}$. Соотношение
сигнал-шум при продолжительности экспозиции $t$ определяется как:
\[
  \mathrm{SNR} \equiv
  \frac{\text{Полное число отсчётов}}{\sqrt{\sum_i \text{Шум}_i^2}}
  = \frac{\text{Полное число отсчётов}}
         {\sqrt{\sigma_{\mathrm{RON}}^2 + \sigma_{\mathrm{DN}}^2 + \cdots}},
\]
\[
  \text{где}\quad
  \sigma_{\mathrm{RON}} = \sqrt{N_{\mathrm{px}} \cdot \mathrm{RON}^2},\quad
  \sigma_{\mathrm{DN}} = \sqrt{N_{\mathrm{px}} \cdot \mathrm{DN} \cdot t}.
\]
